% Caption for cand_c20_combined.png.

\caption{\textbf{Schedule choice sets both task success and actuator energy.}
\textbf{(A)} The curated ladder on the two widowx scenes: six schedules,
30 seeds of 24 episodes each, ordered reference-first then by release period and
labelled with measured period\,/\,observation latency. Boxes are quartiles over
seeds, dots individual seeds, diamonds the mean. Inset frames are the end states
of a successful and a failed rollout. \textbf{(B)} The same measurement opened
out to all 44 CP-SAT operating points and all four workloads: success against
calibrated energy (left), and success on the (period, window) grid, one panel per
workload on its own gradient (right). Stars mark each workload's best schedule by
success. Hatched cells are propagated from a twin whose CP-SAT schedule is
identical to within 0.04\,ms of latency; $\times$ and ? mark infeasible and
unsolved within the solver budget. \emph{Colour means different things in the two
halves}---the arm in (A), the workload in (B)---which is what the A/B split
marks. Energy is calibrated to joules and applied per episode,
$E = c\int\!\sum\tau_{\mathrm{grav}}^2\,dt + P_{\mathrm{idle}}T$, with
$c=0.827\,\mathrm{W/(N\,m)^2}$ and $P_{\mathrm{idle}}=4.61$\,W derived from
ROBOTIS' published stall torque and current for the Dynamixel XM430-W350; the
google constants are a proxy and those series are drawn hollow, since
$P_{\mathrm{idle}}$ dominates their result and calibration does not preserve
their ordering. Plotted energy is the quasi-static lower bound: the raw
$\int\!\sum\tau^2dt$ proxy cannot be converted at face value because the
simulated PD torque runs about 76$\times$ the physical requirement. Derivation in
\texttt{calib/ENERGY\_CALIBRATION.md}.}
